% !TEX program = xelatex
\documentclass[12pt,a4paper]{ctexart}

\usepackage{amsmath,amssymb,amsfonts}
\usepackage{bm}
\usepackage{booktabs}
\usepackage{array}
\usepackage{tabularx}
\usepackage{geometry}
\usepackage{xcolor}
\usepackage{hyperref}
\usepackage{enumitem}
\usepackage{caption}

\geometry{left=2.6cm,right=2.6cm,top=2.6cm,bottom=2.6cm}
\linespread{1.15}
\setlist{nosep,leftmargin=2em}
\captionsetup{font=small,labelfont=bf}
\setlength{\emergencystretch}{3em}

\hypersetup{
  colorlinks=true,
  linkcolor=blue,
  citecolor=blue,
  urlcolor=blue
}

\newcommand{\SPMPC}{SPMPC}
\newcommand{\MPCC}{MPCC}
\newcommand{\MPC}{MPC}
\newcommand{\DWA}{DWA}
\newcommand{\TEB}{TEB}
\newcommand{\DWPP}{DWPP}
\newcommand{\xr}{\bm{x}_{r}}
\newcommand{\xl}{\bm{x}_{l}}
\newcommand{\xs}{\bm{x}_{s}}
\newcommand{\xaug}{\bm{x}}
\newcommand{\vu}{\bm{u}}
\newcommand{\ec}{e_{\mathrm{c}}}
\newcommand{\el}{e_{\mathrm{l}}}
\newcolumntype{Y}{>{\raggedright\arraybackslash}X}

\title{面向移动底盘开口液体运输的在线防晃局部规划\\相关工作组织与方法定位}
\author{}
\date{\today}

\begin{document}
\maketitle

\begin{abstract}
本文档是 \SPMPC{} 论文写作的内部整理稿，用于固定相关工作组织、方法定位、创新边界和实验论证口径。文档重点说明普通移动机器人在线局部规划器、移动底盘液体运输防晃方法以及跨平台液体防晃建模/控制方法之间的层级关系，并明确本文贡献不在于首次将晃液状态引入 \MPC{}，而在于将低阶晃液模态状态嵌入标准轮式移动底盘的在线 \MPCC{} 局部规划层。
\end{abstract}

\tableofcontents
\newpage

\section{定位与创新边界}
\label{sec:positioning}

本文研究面向移动底盘开口液体运输的在线防晃局部规划问题。与单独的液体防晃控制器、离线整段轨迹优化方法或给定轨迹后的跟踪控制不同，本文将低阶晃液状态引入移动机器人在线局部规划层，使规划器在生成短时域底盘运动命令时同时考虑路径跟踪、路径进度、控制平滑性和预测液体响应。

本文的核心主张可概括为：
\begin{quote}
\SPMPC{} 面向标准轮式移动底盘，在滚动时域 \MPCC{} 局部规划问题中显式传播低阶晃液模态状态，并在同一优化问题中联合处理路径误差、路径推进、底盘控制约束、控制平滑性和预测晃液响应，最终输出普通底盘可执行的速度命令。
\end{quote}

本文的创新边界由两个维度共同界定。从平台角度看，本文面向标准轮式移动底盘，目标是输出普通导航栈可直接执行的速度命令；从问题层级角度看，本文处于在线局部规划层，规划器在运行时以滚动时域方式反复求解，而非离线预计算整条轨迹。已有液体防晃方法多数位于离线轨迹优化、给定路径后的速度剖面设计、输入整形、给定轨迹后的跟踪控制或特殊平台控制层；普通移动机器人局部规划器虽然在线可执行，但通常不传播液体动态状态。本文的贡献位于两者的交叉处：在标准轮式移动底盘的在线 \MPCC{} 局部规划问题中显式传播晃液模态状态。

需要指出的是，本文不宣称首次将晃液状态引入 \MPC{} 或预测控制——已有研究已在液体容器转运系统、机械臂、特殊移动平台和给定轨迹跟踪任务中考虑晃液动力学。因此，本文的相关工作和方法表述应围绕“平台”和“问题层级”展开，而不应仅按“是否研究液体晃动”作简单分类。

\section{相关工作定位}
\label{sec:related-work-organization}

液体运输机器人的相关研究大体沿两条主线发展：一类工作显式建模液体晃动，通过轨迹优化、输入整形或预测控制抑制液面振荡；另一类工作关注移动机器人的在线局部规划与实时轨迹生成，强调可执行性、实时性和路径跟踪性能。然而，这两条路线通常并不在同一决策层上相交：前者往往不面向标准轮式移动底盘导航栈中的 online local planner，后者则通常不传播液体动态状态。本文关注的正是这一交叉点。

\subsection{分类原则}
\label{subsec:classification-principle}

相关工作按照两个维度组织。第一个维度是平台是否接近本文任务，包括标准轮式移动底盘、特殊移动平台、机械臂、液罐车、船舶和专用转运机构等。第二个维度是问题层级是否接近本文方法，包括在线局部规划、离线轨迹优化、速度剖面设计、输入整形、跟踪控制、晃液建模和液面测量等。

基于上述原则，本文的相关工作可分为三类。
\begin{enumerate}
  \item \textbf{普通移动机器人在线局部规划器与轨迹优化方法。} 这类方法与本文处于相近的部署层级，能够在线生成底盘可执行命令，但通常不显式传播液体状态。
  \item \textbf{移动底盘液体运输与防晃近邻工作。} 这类工作与本文共享移动平台运输液体的物理任务，但多数属于固定路径、速度剖面、离线轨迹优化、给定轨迹跟踪或特殊平台控制。
  \item \textbf{跨平台液体防晃建模、轨迹生成与预测控制。} 这类工作为低阶晃液建模、输入整形和预测控制提供方法基础，同时也限定了本文创新表述的边界。
\end{enumerate}

这种组织方式的目的不是扩大相关工作范围，而是明确本文与各类已有方法之间的层级关系：普通局部规划方法与本文同层但不具备液体状态预测；移动底盘防晃方法与本文同任务但多数不同层；跨平台防晃方法提供方法基础但通常不构成同层实验基线。

\subsection{普通局部规划器}
\label{subsec:ordinary-local-planners}

移动机器人局部规划器主要解决运行时的短时域运动生成问题，要求在满足运动学、动力学和环境约束的前提下输出可执行底盘命令。经典 \DWA{}、\TEB{} 和 \texttt{mpc\_local\_planner} 分别从速度空间搜索、带时间分配的轨迹优化和滚动时域非线性最优控制等角度构成普通局部规划的重要谱系\cite{D15_Fox1997,D16_Rosmann2017,D17_Rosmann2021}。同时，\MPCC{} 已在移动系统路径进度优化和移动机器人避障局部规划中形成成熟主干\cite{D19_Liniger2015,D20_Brito2019}；RPP、\DWPP{}、LT-DWA、MPPI、差速平台轨迹优化和加加速度受限轨迹生成等已有条目进一步从路径跟踪、动力学可行性、风险处理和平滑性等角度扩展了移动机器人局部规划方法\cite{D01_Macenski2023,D22_Ohnishi2026,D02_Jian2023,D03_Zhang2025,D06_Williams2017,D13_Berscheid2021}。

这些方法与本文具有相同的系统接口特征：它们均服务于移动机器人运行时局部决策，并最终生成底盘可执行命令。因此，它们适合作为候选同层实验基线，用于检验普通局部规划器在液体运输任务中是否已经足够。

然而，这类方法通常只传播机器人位姿、速度、路径误差、障碍距离、控制变化或风险变量，不包含容器内液体的模态位移和模态速度。对于液体运输任务，仅降低控制变化或提高轨迹平滑性并不等价于抑制晃动。两段控制序列可能具有相近的速度、加速度和加加速度特征，但由于与当前液体相位的耦合关系不同，其未来液体响应可能显著不同。因此，普通局部规划器可以作为同层对照，但不能直接表达液体动态记忆。

\subsection{移动底盘防晃}
\label{subsec:mobile-liquid-transport}

移动底盘液体运输是与本文最接近的物理任务。Hamaguchi 及其合作者较早研究了轮式移动机器人或移动台车运输圆柱容器液体时的路径设计、速度剖面、输入整形和轨迹跟踪控制\cite{B04_Hamaguchi2002,B03_Hamaguchi2004,B02_Hamaguchi2005}。这些工作表明，移动平台运动会显著激发容器内液体晃动，运动规划和控制策略需要考虑液体响应。

近期工作进一步将液体模型纳入移动机器人或移动容器的优化与控制。Lim 等将液体晃动近似为球摆动力学，并进行二维轨迹优化以降低晃动\cite{B01_Lim2024}。Nguyen Viet 等将非线性等效质量--弹簧--阻尼模型、时间最优规划、跟踪控制和输出约束结合起来处理扰动下的容器运输问题\cite{B11_Viet2025}。Prabakaran 等在可重构楼梯服务机器人上研究了晃液感知轨迹控制和 \MPC{} 跟踪控制\cite{B12_Prabakaran2026}。

上述工作证明了移动平台液体运输中显式液体模型和优化控制的价值，但它们与本文的主要差异在于问题层级。Hamaguchi 系列工作更偏向给定路径后的速度剖面、输入整形和跟踪控制；Lim 等工作更接近离线整段轨迹优化；Nguyen Viet 等工作采用“参考轨迹规划加跟踪控制”的两层结构；Prabakaran 等工作面向特殊服务机器人平台和给定轨迹控制。相比之下，本文将低阶晃液状态直接并入标准轮式移动底盘的在线 \MPCC{} 局部规划问题，在每个控制周期内重新求解短时域优化并输出底盘速度命令。

\begin{table}[htbp]
  \centering
  \small
  \caption{移动底盘液体运输近邻工作与本文方法的层级差异}
  \begin{tabularx}{\linewidth}{p{0.24\linewidth}p{0.27\linewidth}Y}
    \toprule
    类型 & 代表工作 & 与本文方法的主要差异 \\
    \midrule
    固定路径、速度剖面与输入整形 & Hamaguchi 系列\cite{B04_Hamaguchi2002,B03_Hamaguchi2004,B02_Hamaguchi2005} & 主要在预先设计的路径或速度剖面上抑制晃动，不是在每个控制周期内求解带路径进度的局部 \MPCC{} 问题。 \\
    离线液体约束轨迹优化 & Lim 等\cite{B01_Lim2024} & 显式使用液体动力学模型，但更接近一次性整段轨迹优化，不直接对应运行时局部规划层。 \\
    时间最优规划与跟踪控制 & Nguyen Viet 等\cite{B11_Viet2025} & 以参考轨迹规划和鲁棒跟踪为主要结构，不直接输出标准导航栈局部规划接口下的底盘命令。 \\
    特殊平台晃液感知跟踪控制 & Prabakaran 等\cite{B12_Prabakaran2026} & 说明晃液感知 \MPC{} 或跟踪控制已有研究，但平台、任务接口和控制层级不同于标准轮式移动底盘在线 \MPCC{}。 \\
    \bottomrule
  \end{tabularx}
\end{table}

\subsection{跨平台防晃方法}
\label{subsec:cross-platform-slosh-control}

液体晃动也广泛存在于机械臂、SCARA 机构、液体容器转运系统、液罐车、船舶和专用机构中。这类研究为本文提供了重要的方法基础。例如，低阶晃液估计研究表明，等效摆、质量--弹簧--阻尼模型或模态状态可用于描述主导晃液响应\cite{A02_Guagliumi2021,A01_Leva2021}；输入整形研究表明，命令预整形可降低残余振动\cite{A07_Singer1990}；Okatsuka 等已经在液体容器转运系统中将晃液动力学纳入预测控制框架\cite{A12_Okatsuka2011,C05_Okatsuka2012}；机械臂和机器人操作中的防晃轨迹优化、容器姿态补偿和跟踪控制进一步说明显式液体模型能够改善运动生成与控制性能\cite{C02_Muchacho2022,C03_Moriello2018,C06_Viet2024}。

这些文献的作用主要体现在两个方面。第一，它们支持本文采用低阶晃液状态作为在线规划模型的技术选择。第二，它们说明晃液感知预测控制和防晃轨迹优化已有充分研究，因此本文不能将创新点表述为首次将晃液状态引入 \MPC{}。此外，托盘运输、结构补偿或重构机构等路线也说明移动平台上的防晃并不必然发生在 planner 层；这些工作有助于限定本文的方法边界，但其控制输入、自由度和部署接口不同于标准轮式移动底盘局部规划器，因此通常不作为本文的同层实验基线。

\subsection{小结}
\label{subsec:related-work-summary}

综上，已有研究分别建立了移动机器人在线局部规划、移动底盘液体运输、低阶晃液建模和防晃预测控制的基础，但尚无工作将低阶晃液状态直接嵌入标准轮式移动底盘的在线 \MPCC{} 局部规划层。本文填补这一空白，方法细节见第~\ref{sec:method}~节。

\section{方法概述}
\label{sec:method}

\subsection{问题设定}
\label{subsec:problem-setting}

液体晃动不是底盘加速度的静态函数，而是由历史激励驱动的二阶动态系统。当前模态位移 $\eta$ 和模态速度 $\dot{\eta}$ 决定了未来激励是放大还是抑制晃动。因此，仅在命令层约束平滑性，或事先规划静态速度剖面，均无法在运行时预判液体的动态响应——这正是需要将晃液状态显式引入局部规划预测模型的根本原因。

与“先规划局部路径、再设计速度剖面或跟踪控制器”的两阶段串联方法不同，\SPMPC{} 在单一滚动时域优化中统一处理路径推进、路径误差、底盘约束和晃液响应，输出可直接执行的底盘命令。

本文方法可以概括为一个带液体动态记忆的移动底盘在线局部规划器。它不是单独的液体晃动控制器，也不是离线一次性优化整条轨迹的方法，更不是给定轨迹后的普通跟踪控制器；本文把开口液体运输中的防晃需求放到标准轮式移动底盘的在线局部规划层来处理。具体而言，在给定安全参考路径附近，\SPMPC{} 每个控制周期求解一个同时包含机器人状态、路径进度和低阶晃液状态的短时域 \MPCC{} 问题，并输出普通底盘可执行的速度命令。

该方法不追求高保真自由液面仿真，也不依赖真实液面高度的在线闭环测量。晃液状态在本文中是由低阶模型传播得到的规划内部状态，用来描述液体的主导模态位移、模态速度和当前相位；它的作用是让优化器在预测时域内判断未来底盘激励会放大还是抵消晃动。因此，\SPMPC{} 的核心不是简单增加平滑权重，而是在局部规划器的预测模型中显式加入液体动态状态。

在每个控制周期内，\SPMPC{} 接收四类信息：当前底盘状态，包括位置、航向、线速度和角速度；给定安全参考路径及当前路径进度；上一周期控制量，用于约束控制变化；以及由模型传播得到的当前晃液状态。随后，规划器在同一预测时域内联合使用机器人运动学、路径进度动力学和低阶晃液动力学，优化路径跟踪、路径推进、控制幅值、控制平滑性和预测晃液响应。

滚动时域执行流程为：
\begin{enumerate}
  \item 根据当前底盘状态、路径进度和晃液状态初始化增强状态；
  \item 在预测时域内传播增强状态系统，包括机器人运动学模型、路径进度动力学和低阶晃液模态模型；
  \item 求解包含轮廓误差、滞后误差、路径进度奖励、控制代价、控制差分代价和晃液状态代价的 \MPCC{} 问题；
  \item 只执行第一帧最优控制，并将其转换为标准底盘速度命令；
  \item 下一控制周期重新根据最新状态初始化并求解。
\end{enumerate}

输出为标准差速底盘速度命令，可直接通过移动机器人导航栈的局部规划器接口下发：
\begin{equation}
  \bm{u}_{\mathrm{cmd}} = (v_{\mathrm{cmd}},\,\omega_{\mathrm{cmd}}).
\end{equation}

\subsection{增强状态模型}
\label{subsec:augmented-state-model}

为同时描述路径跟踪、底盘可执行性和路径推进，本文采用包含角速度状态的 \MPCC{} 表述。机器人与路径进度状态定义为
\begin{equation}
  \xr =
  \begin{bmatrix}
    p_x & p_y & \theta & v & s & \omega
  \end{bmatrix}^{\top},
\end{equation}
其中 $(p_x,p_y)$ 为底盘位置，$\theta$ 为航向角，$v$ 为线速度，$s$ 为路径进度，$\omega$ 为角速度。控制输入定义为
\begin{equation}
  \vu =
  \begin{bmatrix}
    a & \alpha & v_s
  \end{bmatrix}^{\top},
\end{equation}
其中 $a$ 为线加速度，$\alpha=\dot{\omega}$ 为角加速度，$v_s=\dot{s}$ 为虚拟路径进度速度。连续时间机器人模型为
\begin{equation}
\begin{aligned}
  \dot{p}_x &= v\cos\theta, &
  \dot{p}_y &= v\sin\theta, &
  \dot{\theta} &= \omega, \\
  \dot{v} &= a, &
  \dot{s} &= v_s, &
  \dot{\omega} &= \alpha .
\end{aligned}
\end{equation}

上述表述使路径进度成为优化变量，而不是将问题简化为给定轨迹后的跟踪控制。同时，将角速度作为状态、角加速度作为控制，可在优化问题内部直接约束转向变化，并为横向晃液激励建模提供状态量。

根据二维平面激励下的低阶晃液估计研究，容器内主导晃液响应可由少量模态坐标及其速度近似描述\cite{A01_Leva2021,A02_Guagliumi2021}。因此，本文不采用高保真自由液面仿真，而是在局部规划层使用控制友好的二阶模态模型。每个水平方向的低阶晃液状态定义为
\begin{equation}
  \bm{x}_{s,i}=\begin{bmatrix}\eta_i & \dot{\eta}_i\end{bmatrix}^{\top},
  \qquad i\in\{x,y\},
\end{equation}
其中 $\eta_i$ 表示主导晃液模态位移，$\dot{\eta}_i$ 表示模态速度。采用二阶模态方程描述晃液响应：
\begin{equation}
  \ddot{\eta}_i + 2\zeta_i\omega_i\dot{\eta}_i + \omega_i^2\eta_i = -\kappa_i a_i,
\end{equation}
其中 $\omega_i$、$\zeta_i$ 和 $\kappa_i$ 分别表示等效固有频率、阻尼比和输入增益，$a_i$ 为容器水平激励。

对于轮式移动底盘，本文采用控制友好的激励近似：
\begin{equation}
  a_x = a, \qquad a_y \approx v\omega .
\end{equation}
该近似表示纵向晃动主要由线加速度激发，横向晃动主要由转弯时的向心加速度激发。

组合机器人状态与液体状态后，增强状态为
\begin{equation}
  \xaug =
  \begin{bmatrix}
    p_x & p_y & \theta & v & s & \omega &
    \eta_x & \dot{\eta}_x & \eta_y & \dot{\eta}_y
  \end{bmatrix}^{\top}.
\end{equation}
该增强状态是本文方法的关键。普通局部规划器只传播机器人运动状态和路径误差，而 \SPMPC{} 进一步传播当前液体相位，使优化器能够区分相同平滑程度但液体响应不同的未来运动。

从状态空间角度看，引入 $\eta_i$ 和 $\dot{\eta}_i$ 的意义在于使局部预测模型包含液体动态记忆，并形成适合有限时域最优控制问题的近似马尔可夫状态描述。若规划器只包含底盘位姿、速度和路径进度，未来液体响应仍依赖未显式表示的历史激励和当前液体相位；加入低阶晃液状态后，离散增强系统可写为
\begin{equation}
  \xaug_{k+1}=f_d(\xaug_k,\vu_k),
\end{equation}
即下一步增强状态由当前增强状态和当前控制决定。这样，液体的动态记忆被显式纳入 OCP 的状态空间，而不是仅通过控制平滑项间接处理。

\subsection{MPCC 优化问题}
\label{subsec:mpcc-ocp}

给定参考路径 $\bm{r}(s) = [r_x(s),\,r_y(s)]^\top$，底盘位置相对于路径的误差向量为
\begin{equation}
  \bm{e}(s) = [p_x - r_x(s),\; p_y - r_y(s)]^\top.
\end{equation}
路径切向 $\bm{t}(s) = \bm{r}'(s)/\|\bm{r}'(s)\|$，法向 $\bm{n}(s) = \bm{t}(s)^\perp$。轮廓误差与滞后误差定义为
\begin{equation}
  \ec = \bm{n}(s)^{\top}\bm{e}(s),
  \qquad
  \el = \bm{t}(s)^{\top}\bm{e}(s).
\end{equation}
其中 $\ec$ 衡量横向偏离程度，$\el$ 描述沿路径方向的进度误差。

离散化增强动力学后，每个控制周期求解如下有限时域最优控制问题：
\begin{equation}
\begin{aligned}
  \min_{\{\xaug_k,\vu_k\}_{k=0}^{N-1}} \quad
  & \sum_{k=0}^{N-1}
  \Big(
    q_c e_{\mathrm{c},k}^{2}
    + q_l e_{\mathrm{l},k}^{2}
    - q_p v_{s,k}
    + \|\vu_k\|_{R}^{2}
    + \|\vu_k-\vu_{k-1}\|_{S}^{2}
    + \|\bm{x}_{s,k}\|_{Q_s}^{2}
  \Big) \\
  & + q_{c,N}e_{\mathrm{c},N}^{2}
    + q_{l,N}e_{\mathrm{l},N}^{2}
    + \|\bm{x}_{s,N}\|_{Q_{s,N}}^{2} \\
  \mathrm{s.t.}\quad
  & \xaug_0=\hat{\xaug}(t), \\
  & \xaug_{k+1}=f_d(\xaug_k,\vu_k),\qquad k=0,\ldots,N-1, \\
  & \xaug_k\in\mathcal{X},\qquad \vu_k\in\mathcal{U} .
\end{aligned}
\end{equation}

目标函数中，轮廓误差和滞后误差保证路径跟踪性能，路径进度奖励避免规划器通过停滞来降低晃动，控制幅值和控制差分约束保证底盘命令的可执行性和连续性，晃液状态代价则使优化器在预测时域内主动抑制液体响应。

\begin{table}[htbp]
  \centering
  \small
  \caption{\SPMPC{} 目标函数各项的作用}
  \begin{tabularx}{\linewidth}{p{0.25\linewidth}p{0.36\linewidth}Y}
    \toprule
    代价项 & 优化作用 & 在论文论证中的作用 \\
    \midrule
    $q_c\ec^2$, $q_l\el^2$ & 控制横向路径误差和沿路径进度误差 & 保持局部规划器的路径跟踪能力。 \\
    $-q_p v_s$ & 奖励沿参考路径推进 & 避免规划器仅通过降低速度或停滞来抑制晃动。 \\
    $\|\vu\|_R^2$ & 限制控制幅值 & 保证底盘命令满足执行约束。 \\
    $\|\vu_k-\vu_{k-1}\|_S^2$ & 平滑控制变化 & 保留普通平滑规划器的基本优势。 \\
    $\|\bm{x}_{s}\|_{Q_s}^2$ & 惩罚预测晃液状态 & 引入液体动态状态，使规划器具备晃液感知能力。 \\
    \bottomrule
  \end{tabularx}
\end{table}

典型约束包括底盘速度、角速度、线加速度、角加速度和路径进度速度约束：
\begin{align}
  0 \leq v_k \leq v_{\max}, \qquad
  |\omega_k| \leq \omega_{\max}, \\
  |a_k| \leq a_{\max}, \qquad
  |\alpha_k| \leq \alpha_{\max}, \qquad
  0 \leq v_{s,k} \leq v_{s,\max} .
\end{align}

在需要时，也可加入模型预测晃液软约束：
\begin{equation}
  \eta_{x,k}^{2}+\eta_{y,k}^{2} \leq \eta_{\max}^{2}+\epsilon_k,
  \qquad \epsilon_k\geq 0.
\end{equation}
该约束仅表示基于低阶模型的运行边界，不构成严格的无溢出保证。尽管如此，软约束通过在预测时域内引入边界惩罚区，使 solver 在可行时倾向于生成远离晃液上限的轨迹，从而在统计意义上降低液面越限概率；与硬约束相比，软约束在激励较强时不会导致 OCP 不可行，有助于维持求解鲁棒性。

求解完成后，控制器仅执行第一帧优化结果：
\begin{equation}
  v_{\mathrm{cmd}} = v_0+a_0^{\star}\Delta t,
  \qquad
  \omega_{\mathrm{cmd}}=\omega_0+\alpha_0^{\star}\Delta t.
\end{equation}
随后系统进入下一控制周期，并重新根据当前状态求解滚动时域优化问题。


\subsection{消融设计}
\label{subsec:ablation-design}

降低控制激励有助于减少晃动：经典输入整形\cite{A07_Singer1990}和加加速度受限轨迹生成\cite{D13_Berscheid2021}均表明，对命令平滑处理可降低残余振动。然而，平滑性是防晃运动的必要条件，但不是充分条件。

液体响应具有动态记忆——当前模态位移 $\eta_i$ 和模态速度 $\dot\eta_i$ 决定了未来激励是放大还是抵消晃动。因此，两段在速度、加速度或控制差分上同样平滑的运动，在液体已有初始相位（$\eta\neq 0$）的情况下，仍可能产生截然不同的液体响应。这种“相位耦合”效应在连续机动过程中尤为显著：上一次转弯留下的模态振荡会与下一次激励叠加或抵消，仅靠控制平滑性无法预判结果。

需要说明的是，当液体初始静止（$\eta=0,\,\dot\eta=0$）时，显式晃液状态预测与平滑基线的区别缩小——两者均从零初始状态出发，优化目标本质相近。因此，本文的消融实验将重点考察连续机动场景（液体存在初始相位），以检验显式晃液状态预测是否具有超出单纯控制平滑性的贡献。后续实验中，\texttt{B\_smooth} 与 \texttt{B\_ours} 的对比用于系统量化这一差异。

\medskip
\noindent\textbf{核心结论：同样平滑的控制序列，不意味着同样低的液面响应。}

为区分晃液状态预测和平滑性对性能的贡献，后续实验采用四类内部变体：
\begin{table}[htbp]
  \centering
  \small
  \caption{内部消融变体及其验证目的}
  \begin{tabularx}{\linewidth}{p{0.18\linewidth}p{0.22\linewidth}p{0.22\linewidth}Y}
    \toprule
    变体 & 晃液模型与代价 & 平滑权重 & 验证目的 \\
    \midrule
    \texttt{B0} & 不启用 & 弱 & 检验基础 \MPCC{} 是否能够完成路径跟踪。 \\
    \texttt{B\_smooth} & 不启用 & 强 & 检验仅增强控制平滑性是否足以降低晃动。 \\
    \texttt{B\_slosh} & 启用 & 弱 & 检验显式晃液状态预测是否具有独立贡献。 \\
    \texttt{B\_ours} & 启用 & 强 & 检验完整方法是否能够同时兼顾跟踪、平滑性和防晃性能。 \\
    \bottomrule
  \end{tabularx}
\end{table}

\section{实验设计}
\label{sec:experiment-rationale}

当前实验设计应服务于三类问题。第一，普通局部规划器是否已经足够处理液体运输任务；第二，已有移动底盘防晃方法与本文方法在层级和接口上有何差异；第三，本文性能提升是否来自显式晃液状态预测，而不是单纯降低速度或增强平滑性。

基于这一目标，后续实验可优先考虑以 \DWA{}、\TEB{}、\texttt{mpc\_local\_planner} 和 LT-DWA 作为候选同层基线，用于代表不具备液体状态预测的普通局部规划器；RPP、\DWPP{}、MPPI 和加加速度受限轨迹生成等方法可作为普通局部规划谱系中的写作参照。Lim 和 Hamaguchi 系列可作为同物理任务但不同规划层级的近邻对比；\texttt{B0}、\texttt{B\_smooth}、\texttt{B\_slosh}、\texttt{B\_ours} 及 hard-cap 变体构成内部消融。固定路径仿真实验可以作为第一版论文的主要验证场景，点到点导航、复杂障碍物避让和实物实验可作为后续系统扩展。

\begin{table}[htbp]
  \centering
  \small
  \caption{实验对比对象与论证作用}
  \begin{tabularx}{\linewidth}{p{0.22\linewidth}p{0.31\linewidth}Y}
    \toprule
    对比对象 & 代表方法 & 论证目的 \\
    \midrule
    普通局部规划器 & \DWA{}、\TEB{}、\texttt{mpc\_local\_planner}、LT-DWA & 检验无液体状态预测的同层 planner 是否足以处理开口液体运输任务。 \\
    普通规划谱系参照 & RPP、\DWPP{}、MPPI、Ruckig & 支撑 related work 中“同层但不传播液体状态”的方法谱系，不一定全部作为第一版实验基线。 \\
    移动底盘防晃近邻 & Hamaguchi 系列、Lim 等 & 说明同物理任务已有防晃规划/控制研究，但其规划层级和接口不同于本文。 \\
    内部消融 & \texttt{B0}、\texttt{B\_smooth}、\texttt{B\_slosh}、\texttt{B\_ours}、hard-cap 变体 & 区分路径跟踪、控制平滑性、显式晃液状态预测和晃液边界约束各自的贡献。 \\
    \bottomrule
  \end{tabularx}
\end{table}

\subsection{仿真结果}
\label{subsec:simulation-result-summary}

在固定 S 曲线液体运输仿真中，当前结果采用 strict fresh-sim、$N=9$ 的到点口径，并以 slosh peak、slosh p95 和 slosh RMS 作为液体响应指标。结果表明，加入晃液感知代价或边界后，\SPMPC{} 在保持可达性的同时显著降低了高分位晃动趋势。

以 1\,mm 评价口径的内部消融为例，\texttt{B\_slosh} 将 full-run slosh p95 从 \texttt{B0} 的 0.594\,mm 降至 0.496\,mm，降低 16.5\%；core p95 从 0.540\,mm 降至 0.413\,mm，降低 23.5\%。\texttt{B\_ours} 更适合说明峰值抑制：full-run peak 从 1.268\,mm 降至 1.012\,mm，降低 20.2\%；core peak 从 1.256\,mm 降至 0.960\,mm，降低 23.6\%。加入 hard-cap 后，\texttt{B\_ours\_hard} 将 core p95 进一步降至 0.390\,mm，将 event p95 降至 0.452\,mm，分别较 \texttt{B0} 降低 27.8\% 和 18.1\%。

在外部 fixed-path $N=9$ 对比中，\texttt{spmpc\_B\_ours\_hard\_1mm} 的 slosh peak、p95 和 RMS 分别为 1.062\,mm、0.473\,mm 和 0.257\,mm，均为 \DWA{}、\TEB{}、\texttt{mpc\_local\_planner} tuned、LT-DWA、Hamaguchi profile 和 Lim profile 等对比方法中的最低值。相对 \DWA{}、\TEB{}、\texttt{mpc\_local\_planner} tuned 和 LT-DWA，其 p95 分别降低 70.9\%、71.3\%、72.0\% 和 73.8\%；相对 Hamaguchi profile 和 Lim profile，其 p95 分别降低 56.2\% 和 62.8\%。这些结果支持本文的实验主张：显式液体状态预测和晃液边界约束主要压低的是高分位振荡尾部，而不仅是让底盘命令更平滑。

\begin{table}[htbp]
  \centering
  \small
  \caption{固定 S 曲线仿真中的晃液指标摘要}
  \begin{tabularx}{\linewidth}{p{0.25\linewidth}p{0.30\linewidth}p{0.16\linewidth}Y}
    \toprule
    结果类型 & 主要数值 & 改善幅度 & 论证作用 \\
    \midrule
    \texttt{B\_slosh} vs. \texttt{B0} & full p95: 0.594 $\rightarrow$ 0.496\,mm；core p95: 0.540 $\rightarrow$ 0.413\,mm & 16.5\%；23.5\% & 说明显式晃液状态预测能够降低高分位趋势。 \\
    \texttt{B\_ours} vs. \texttt{B0} & full peak: 1.268 $\rightarrow$ 1.012\,mm；core peak: 1.256 $\rightarrow$ 0.960\,mm & 20.2\%；23.6\% & 说明完整方法能够抑制峰值响应。 \\
    \texttt{B\_ours\_hard} vs. \texttt{B0} & core p95: 0.540 $\rightarrow$ 0.390\,mm；event p95: 0.552 $\rightarrow$ 0.452\,mm & 27.8\%；18.1\% & 说明 hard-cap 主要压低高曲率窗口的尾部风险。 \\
    外部同层 baseline & peak/p95/RMS: 1.062/0.473/0.257\,mm & p95 低 70.9--73.8\% & 相对 \DWA{}、\TEB{}、\texttt{mpc\_local\_planner} 和 LT-DWA 均取得最低晃液指标。 \\
    profile baseline & p95: 0.473\,mm & p95 低 56.2\%/62.8\% & 相对 Hamaguchi/Lim profile 说明在线晃液预测具有额外价值。 \\
    \bottomrule
  \end{tabularx}
\end{table}

实验结果应同时报告晃液指标、路径跟踪指标、任务效率指标和实时性指标。若当前阶段以仿真为主，还应通过模型参数失配或初始晃液状态扰动检验方法的稳健性，避免结果仅反映“用同一低阶模型优化同一低阶指标”的情况。

\textbf{关于当前实物实验状态：}初步实物录包（2026-07-02）发现仿真优势在实物上尚未完全复现，主要原因已定位为：硬约束参数命名空间错误导致实际约束未生效；延迟补偿参数偏差（角速度配置220\,ms，实测约5--10\,ms）可能导致相位补偿方向错误。上述问题均在代码层面识别并已修复，延迟值校准实验正在准备中。因此，当前不宜将初步实物结果解释为“方法本身无效”，应在参数落地确认和延迟校准完成后重跑验证。

\section{后续事项}
\label{sec:todo}

\begin{enumerate}
  \item \textbf{相关工作写作：} “平台$\times$问题层级”分类框架已确定，后续需在各段落中落实对代表文献的逐篇定位陈述；\DWA{}、\TEB{}、\texttt{mpc\_local\_planner}、移动机器人 \MPCC{} 主干文献和局部规划评价基准（MRPB）应作为主文必引。
  \item \textbf{近邻工作定位：} Hamaguchi 系列和 Lim 等适合作为同物理任务但不同层级的概念对比；是否纳入实验 baseline 取决于公开实现的可获取性。
  \item \textbf{方法表述：} 创新点表述已统一为“标准轮式移动底盘在线 \MPCC{} 局部规划层中的晃液状态增强”，后续写作避免使用“首次将晃液状态引入 \MPC{}”的措辞。
  \item \textbf{模型边界：} 当前低阶二阶模态模型已足以支撑第一版论文；正文需明确说明模型预测晃液状态是规划器内部代理量，不等同于真实液面高度观测。
  \item \textbf{实验验证：} \texttt{B\_ours} 与 \texttt{B\_smooth} 的对比是回应“性能提升来自平滑性还是晃液感知”的核心实验证据，应在实验章节中优先呈现。
\end{enumerate}

\section{汇报表述}
\label{sec:short-summary}

本文可概括为以下表述：
\begin{quote}
普通移动机器人局部规划器能够在线生成底盘可执行命令，但通常不传播液体动态状态；已有液体防晃方法显式考虑晃动，但多数位于离线轨迹优化、给定路径速度剖面、输入整形、给定轨迹跟踪或特殊平台控制层。本文面向标准轮式移动底盘，将低阶晃液模态状态嵌入在线 \MPCC{} 局部规划问题，在滚动时域内联合优化路径跟踪、路径进度、底盘控制平滑性和预测液体响应，并输出标准速度命令。
\end{quote}

\bibliographystyle{IEEEtran}
\bibliography{../spmpc_paper_cn/references}

\end{document}
